\section{Conclusions}
\label{sec:conclusions}

This project set out to test whether machine learning can improve
polygenic risk prediction for autism spectrum disorder beyond the
canonical pruning-and-thresholding (P+T) and Bayesian-shrinkage
(PRS-CS) baselines. Working purely from two public summary-statistics
releases (the iPSYCH--PGC ASD November~2017 dataset and the Stein-Lab
SPARK~+~iPSYCH~+~PGC meta-analysis), we built an end-to-end pipeline
that harmonises the two files across genome builds, derives an
independent SPARK-only target per SNP, and benchmarks four SNP-level
scoring methods on a chromosomal hold-out.

\subsection{Headline findings}

\begin{itemize}[leftmargin=1.4em]
  \item The two public sumstats files use different genome builds.
        Joining on rsID instead of chromosome and position recovers
        \textbf{roughly 100 times more} overlapping SNPs (about
        7.4\,million vs.\ 73\,thousand) and yields \textbf{6\,309\,737}
        SNPs after quality control.
  \item The SPARK-only back-calculation described in the Methods
        passes an internal consistency check (rebuilding the published
        meta Z to machine precision), so we can train and test without
        individual-level genotypes.
  \item On the held-out fold (chromosomes 21--22, 170\,039 SNPs), the
        \textbf{XGBoost re-weighter ranks first} with correlation
        $r = 0.058$ (95\,\% CI 0.053--0.063). The neural network is
        next ($r = 0.048$), then empirical-Bayes shrinkage ($r =
        0.041$), then the discovery-only GNN fallback ($r = 0.037$).
        The bootstrap intervals for these four do not overlap, so the
        ordering is statistically clear.
  \item P+T underperforms at every threshold we tried. Even when almost
        all SNPs are nominally kept ($p \leq 1$), only 213 test
        SNPs remain after clumping on chromosomes 21--22; at the usual
        genome-wide significance cut-off, none remain. P+T is a poor
        fit for a highly polygenic trait like ASD.
  \item For calibration, the MLP is closest to predicting the right
        scale of Z (slope about 0.9 vs.\ an ideal 1.0); XGBoost is
        somewhat too aggressive; empirical Bayes is far too conservative.
\end{itemize}

\subsection{What the results mean}

The main positive result is that an ML re-weighter can beat a standard
Bayesian-style baseline using only the discovery Z-score, its absolute
value, imputation quality, and simple summaries of neighbouring SNPs
in a 500\,kb window. Those neighbourhood summaries let the model
up-weight regions where many SNPs look interesting together -- a kind
of local smoothing that a single global shrinkage rule cannot mimic.
We did not yet load gene annotations (Gencode, SFARI, STRING, etc.); once
those are added, we expect further gains.

The main caution for P+T is visible in the table: throwing away
everything except ``genome-wide significant'' SNPs leaves almost
nothing to score on small chromosomes, which is exactly why many
studies have moved to continuous shrinkage or learned weights.

\subsection{Limitations}

\begin{itemize}[leftmargin=1.4em]
  \item \textbf{Modest raw correlations.} Reported $r$ values around
        0.04--0.06 are small in absolute terms because the SPARK-only
        target is estimated from a smaller effective sample than the
        full meta-analysis, so each SNP is very noisy. The \emph{order}
        of methods is still meaningful for comparing approaches.
  \item \textbf{Missing annotation files.} Without gene coordinates
        and pathway databases locally, the gene-network path could not
        map SNPs to genes, and several feature columns were all zeros.
        The ML advantage we see is therefore a lower bound on what the
        full pipeline could do.
  \item \textbf{No individual-level evaluation.} We score SNP-level
        agreement, not whether the final PRS separates cases and controls
        in a real cohort. That step needs individual genotypes and
        phenotypes (e.g.\ SPARK with appropriate approval).
  \item \textbf{Single ancestry.} Both releases are predominantly
        European-ancestry; results may not transfer to other
        populations without retraining or new references.
\end{itemize}

\subsection{Future work}

The next steps are practical: download the reference panels and
annotation files the Makefile lists, re-run so gene mapping and biology
features are live, optionally swap in the full PRS-CS program with its
LD reference, and validate on individual-level data when access is
available. The code is already organised so each of those is a separate
make target.
